\documentclass[11pt,a4paper]{article}
\usepackage{fontspec}
\usepackage{amsmath,amssymb,amsthm}
% 0.9in rather than Paper C's 1in: this note is mostly wide numeric tables, and the extra
% 14.4pt of text width is what Table 1 needs at the 2^51 floor.
\usepackage[margin=0.9in]{geometry}
\usepackage{longtable,booktabs,array,calc}
\usepackage{graphicx}
\usepackage{float}
\floatplacement{figure}{H}
\usepackage{needspace}
\allowdisplaybreaks[1]
\clubpenalty=10000
\widowpenalty=10000
\raggedbottom
\usepackage{fvextra}
\DefineVerbatimEnvironment{verbatim}{Verbatim}{breaklines=true,fontsize=\small}
\usepackage{xurl}
\usepackage[hidelinks]{hyperref}
\hypersetup{pdftitle={No m-cycles of the 3n−1 map for m ≤
58},pdfauthor={Philippe Cochin},pdfsubject={Mathematics preprint}}
% The Markdown keeps "≤" in its title and in one reference; Latin Modern's text
% faces lack the glyph, so it is typeset as \le. Set after \hypersetup so the
% PDF title keeps the character itself.
\catcode`\≤=\active
\def≤{\ensuremath{\le}}
\urlstyle{tt}
\setlength{\emergencystretch}{3em}
\setlength{\LTpre}{1em}
\setlength{\LTpost}{1em}
% Table 1 carries six columns of up to sixteen digits at the 2^51 floor; at 3pt its twelve
% gaps put it 9.5pt over the text width.
% Table 1 at the 2^51 floor: six columns, ceilings of up to sixteen digits. The filter sets
% tables \footnotesize; with the 0.9in margins above, 2pt gutters leave a few points of slack.
\setlength{\tabcolsep}{2pt}
\renewcommand{\arraystretch}{1.12}
\providecommand{\tightlist}{\setlength{\itemsep}{0pt}\setlength{\parskip}{0pt}}
\providecommand{\passthrough}[1]{#1}
\setcounter{secnumdepth}{-1}
\title{No m-cycles of the 3n−1 map for m ≤ 58\\[0.5em]{\large The
Simons--de Weger template on the negative side, from a verification
floor of \(2^{51}\)}}
\author{Philippe Cochin \\ \texttt{philippe@cochin.fr} \\
  \small ORCID \href{https://orcid.org/0009-0004-1939-3382}{0009-0004-1939-3382}}
\date{21 September 2026}
\begin{document}
\maketitle
\textbf{Preprint.} One external input, Rhin's effective measure for
linear forms in \(\log2\) and \(\log3\); one computational input, the
laboratory's own verification floor, a CPU certificate to \(2^{44}\) and
a GPU sweep to \(2^{51}\). Everything else is proved below or is exact
arithmetic recorded with its margins, and every number in the tables is
recomputed by a second, independent route (Section 8). Neither the
\(3n-1\) cycle question nor divergence on either map is settled here.

\begin{abstract}

Let \(g(y)=y/2\) for even \(y\) and \(g(y)=(3y-1)/2\) for odd \(y\), the
\(3n-1\) map on the positive integers, equivalently the shortcut
\(3n+1\) map on the negative integers. Its known cycles are \(1\),
\((5,7,10)\) and the eleven-element cycle at \(17\), and every start
below \(2^{51}\) reaches one of them. An \(m\)-cycle is a cycle with
\(m\) maximal runs of odd elements. We transpose the template of Simons
and de Weger for \(3n+1\) \(m\)-cycles to this map, deriving its
constants on this side rather than borrowing them: the odd step
subtracts, so in the variable \(u=y-1\) an odd run is exact
multiplication by \(3/2\), a run of \(a\) odd steps starts at
\(y\ge2^a+1\), the cycle equation bounds the linear form
\(\Lambda=o\,\log 3-K\,\log 2\) by \(m/(x_{\min}-1)\) with constant one,
and successive local minima obey \(u_{i+1}<u_i^{\log_2 3}/2\). With
Rhin's bound this gives: the \(3n-1\) map has no \(m\)-cycle with
\(1\le m\le58\) other than the two known ones. For \(m\le2\) that is a
floor-dependent form of a theorem Simons proved without any floor; the
statement is new for \(3\le m\le58\). For \(m\le52\) no admissible cycle
length lies below Rhin's ceiling; for \(53\le m\le58\) the admissible
lengths are excluded by the chaining, the closest by \(9.8\) bits. At
\(m=59\) two lengths remain; they are the output of the template's last
step, not its input, and the floors that remove them are \(2^{51.87}\)
and \(2^{55.18}\). The same tables give \(m\le44\) from \(2^{40}\),
\(m\le49\) from \(2^{44}\), \(m\le63\) from \(2^{56}\), \(m\le68\) from
\(2^{60}\) and \(m\le82\) from \(2^{68}\); on the \(3n+1\) side at the
floor of Simons and de Weger the same enumeration and tests return their
Lemma 18 to the unit. No published verification floor and no \(m\)-cycle
theorem with \(m\ge3\) is known to us for this map.

\end{abstract}

\textbf{2020 Mathematics Subject Classification:} Primary 11B83;
Secondary 11J86, 11Y16.

\textbf{Keywords:} 3n−1 map; Collatz cycles; m-cycles; linear forms in
logarithms; verification floor.

\section{1. The map, the cycles, and what is
known}\label{the-map-the-cycles-and-what-is-known}

Write \(g\) for the shortcut \(3n-1\) map above. Its orbits on the
positive integers are the orbits of the shortcut Collatz map
\(T(x)=x/2\), \((3x+1)/2\) on the negative integers under \(x=-y\):
\((3(-y)+1)/2=-(3y-1)/2\). Three cycles are known, \[
(1),\qquad(5,7,10),\qquad(17,25,37,55,82,41,61,91,136,68,34),
\] with \((K,o)=(1,1),(3,2),(11,7)\) steps and odd steps. Whether these
are all the cycles, and whether every orbit is bounded, are the \(3x-1\)
questions of Lagarias's bibliographies {[}L85{]}, open.

\textbf{The floor.} Every \(1\le y<2^{51}\) reaches one of the three
cycles. Below \(2^{44}\) this is the laboratory's CPU computation of
19--20 September 2026, descent by strong induction on odd starts in one
hundred and twelve chunks: sixteen with a plain walker through
\(2^{40}\) (greatest step count \(544\) against a cap of \(4000\)), then
ninety-six chunks of \(5\cdot2^{35}\) starts with a residue sieve modulo
\(2^{24}\) and a \(2^{16}\)-step block map as an accelerator (the sieve
is exact by Lemma 5), \(8246337208320\) odd starts counted exactly
(greatest step count \(704\) against a cap of \(40000\), counted in
blocks of sixteen); the seven odd starts the first sixteen chunks had
skipped, found from their own printed counts, are verified by full
iteration. The range \([2^{44},2^{51})\) was verified on 21 September
2026 by a CUDA implementation of the same descent on one RTX 5090: the
same sieve, then Barina's domain switch {[}B21{]} with the sign flipped,
which is Lemma 1 below, a whole odd run in one multiplication and a
whole even run in one shift, on a 128-bit state with an overflow report.
It was calibrated first on the CPU's range, where it reproduces every
count, both peaks and the plain walker's step counts exactly, and
cross-checked against the CPU walker on three \(2^{37}\) windows inside
the new range; \(1117103813820416\) odd starts, no failure, no new
cycle, no overflow, greatest step count \(847\), peak about
\(2^{97.0}\), in \(3517\) seconds. No start failed and no new cycle
appeared in either range. The verifier sources, the chunk reports, the
calibration and the coverage checks are archived with the branch
\texttt{\nolinkurl{negative_floor_3x1}}. We found no published floor for
this map: the only statement in print that the searches of Section 8
reached is an unattributed comment on OEIS A037084, ``up to at least
100000000''. The certificate is a computation and is labelled as one
throughout.

\textbf{\(m\)-cycles.} A cycle with at least one even step splits into
\(m\ge1\) maximal runs of odd elements alternating with runs of even
elements. The first element of an odd run is a local minimum, the
element after its last odd step a local maximum. The two known cycles
with an even step have \(m=1\) and \(m=2\). Simons {[}S07{]} proves that
\(g\) has exactly one nontrivial \(2\)-cycle, at \(17\), and does so
without any floor, by de Weger's bounds on \(0<3^K-2^{K+L}<3^{0.89K}\),
Steiner's \(1\)-cycle theorem {[}St77{]} and an extremal inequality; he
states that the floor-free route stops at \(m\ge3\). For \(3n+1\),
Simons {[}Si05{]} excludes \(2\)-cycles, Simons and de Weger {[}SdW{]}
exclude \(m\)-cycles for \(m\le68\) from the floor of November 2004,
\(X_0=301\cdot2^{50}>3.3889\cdot10^{17}\), and Hercher {[}H23{]} reaches
\(m\le91\) from \(695\cdot2^{60}\) with the arrangement of the valleys;
Sinisalo {[}Si03{]} tabulates the cycle lengths a floor alone leaves on
that side, the table this note's \(m\)-free row transposes. The template
of {[}SdW{]} has five steps: a floor; a bound on \(\Lambda\) from the
cycle equation (their Lemma 4 and Corollary 5); a chaining of the local
minima (Lemmas 6, 7); Crandall's lemma generalized for a lower bound on
the length (Lemma 10, Corollary 11) against Rhin's bound for an upper
one (Lemmas 12, 14), sharpened through the partial quotients of
\(\log_23\) (Lemma 16); and de Weger's approximation lattice where a
window remains (Lemma 18). At their floor the three stages give
\(m\le57\), \(m\le63\) and \(m\le68\). Nobody has run any of it on
\(3n-1\), because the first step was missing. {[}S08{]} carries the
template to \(3x+q\) and to \(px+q\), but under the standing hypothesis
\(q=1\) or \(q\ge5\) prime, which excludes \(q=-1\); its Section 5
treats Guy's permutation, a different map, and the paper makes no
statement about \(3x-1\) anywhere. This note runs all five steps with
the constants of this side. The exact enumeration of the admissible
lengths by the three-gap walk is the list their approximation lattice
produces ({[}SdW{]} Section 7: a two-dimensional lattice whose points of
small norm are the pairs \((K,L)\) in the window of Corollary 5 below
the ceiling, found by a reduced basis and a search, each then ``checked
for fulfilling Corollary 5 and Lemma 7''); the two tests applied here to
each length are those two. Run on the \(3n+1\) side at their floor, the
machinery below returns their Lemma 18 (Section 5).

\section{2. Notation}\label{notation}

Throughout \(x=\log2/\log3\), \(\delta=\log_23=1/x\). For a cycle \(C\)
of \(g\) with \(K\) steps and \(o\) odd steps put \[
\Lambda(C)=o\log3-K\log2 .
\] For odd \(y\) put \(u=y-1\). The local minima of an \(m\)-cycle are
\(y_1,\dots,y_m\) in cyclic order, \(u_i=y_i-1\), the run from \(y_i\)
has \(a_i\) odd steps, and \(r_i\ge1\) halvings follow it. \(x_{\min}\)
is the least element of \(C\). The floor is \(X_0=2^{51}\).

\section{3. Five lemmas on the negative
side}\label{five-lemmas-on-the-negative-side}

\textbf{Lemma 1 (odd runs).} Let \(y\ge3\) be odd, \(u=y-1\),
\(a=v_2(u)\). Then \[
g^k(y)-1=\Bigl(\tfrac32\Bigr)^k(y-1)\qquad(0\le k\le a),
\] \(g^k(y)\) is odd for \(k<a\) and \(g^a(y)\) is even. Consequently a
run of \(a\) odd steps starts at \(y\ge2^a+1\), and its local maximum
\(M=g^a(y)\) satisfies \(M\le u^{\delta}+1\).

\emph{Proof.} For odd \(y\), \(g(y)-1=(3y-1)/2-1=\tfrac32(y-1)\). While
the iterates stay odd this repeats, so \(g^k(y)=1+3^ku/2^k\) as long as
\(g^j(y)\) is odd for \(j<k\); and \(1+3^ku/2^k\) is an odd integer
exactly when \(2^{k+1}\mid u\), since \(3^k\) is odd. Hence the first
\(a=v_2(u)\) iterates are odd and the next is even. Then \(u\ge2^a\)
gives \(y\ge2^a+1\), and
\((3/2)^a\le(3/2)^{\log_2u}=u^{\log_2(3/2)}=u^{\delta-1}\) gives
\(M=1+(3/2)^au\le1+u^{\delta}\). \(\square\)

This is Hercher's Lemma 8 {[}H23{]} with \(y\equiv1\pmod{2^a}\) in place
of \(y\equiv-1\): the floor \(2^a+1\) is larger by two, and it is
attained at \(5\) (\(a=2\)) and \(17\) (\(a=4\)).

\textbf{Lemma 2 (cycle equation).} On a cycle \(C\) with at least one
even step, \[
\Lambda(C)=-\sum_{y\in C\ \mathrm{odd}}\log\Bigl(1-\frac1{3y}\Bigr)>0,
\qquad\text{so}\qquad 3^o>2^K,
\] and \(\Lambda(C)<\sum_{y\in C\ \mathrm{odd}}\dfrac1{3y-1}\). The run
from a local minimum \(y\) contributes less than \(1/(y-1)\). Hence, for
an \(m\)-cycle, \[
\Lambda(C)<\sum_{i=1}^m\frac1{y_i-1}\le\frac m{x_{\min}-1}.
\]

\emph{Proof.} An odd step multiplies by \(\tfrac32(1-\tfrac1{3y})\) and
an even step by \(\tfrac12\); around the cycle the product is \(1\), so
\(3^o2^{-K}\prod_{\mathrm{odd}}(1-1/(3y))=1\), which is the displayed
identity, and the product is below \(1\). The bound
\(-\log(1-t)\le t/(1-t)\) at \(t=1/(3y)\) gives \(1/(3y-1)\). Along the
run from \(y\), Lemma 1 gives \(3g^k(y)-1=3(3/2)^ku+2>3(3/2)^ku\), so
the run contributes less than \(\tfrac1{3u}\sum_{k\ge0}(2/3)^k=1/u\).
Finally the least element of \(C\) is odd (an even least element would
halve to something smaller) and its predecessor is even (an odd
predecessor \(y'\) would have \(y'<(3y'-1)/2=x_{\min}\)), so
\(x_{\min}\) is a local minimum and every \(y_i\ge x_{\min}\).
\(\square\)

This is {[}SdW{]} Lemma 4 and Corollary 5 with the sign flipped. The
corresponding bound on the positive side, Hercher's Corollary 17,
carries the constant \(97/54\), and {[}SdW{]} Lemma 4 has \(\sum1/x_i\)
over the minima themselves. Here it is \(1/(y_i-1)\): the odd step
subtracts, so the run's elements sit at or above the geometric
progression rather than below it.

\textbf{Lemma 3 (chaining).} For consecutive local minima,
\(u_{i+1}<u_i^{\delta}/2\). Hence, with
\(B(m)=(\delta^m-1)/(\delta-1)\), \[
o=\sum_{i=1}^ma_i\ \le\ B(m)\log_2u_1-\frac{B(m)-m}{\delta-1},
\qquad K<\delta\,o,
\] and therefore, taking \(y_1=x_{\min}\), \[
\log_2(x_{\min}-1)\ >\ L_{\min}(K,m):=\frac1{B(m)}\Bigl(\frac K\delta+\frac{B(m)-m}{\delta-1}\Bigr).
\]

\emph{Proof.} \(y_{i+1}=M_i/2^{r_i}\) with \(r_i\ge1\) and
\(M_i\le u_i^{\delta}+1\) (Lemma 1), so
\(u_{i+1}=M_i/2^{r_i}-1\le(u_i^{\delta}+1)/2-1<u_i^{\delta}/2\). Taking
logarithms, \(\log_2u_{i+1}<\delta\log_2u_i-1\), and by induction
\(\log_2u_i\le\delta^{i-1}\log_2u_1-(\delta^{i-1}-1)/(\delta-1)\).
Summing over \(i\) and using \(a_i=v_2(u_i)\le\log_2u_i\) gives the
bound on \(o\); \(K\log2<o\log3\) is \(\Lambda>0\). Solving for
\(\log_2u_1\) gives \(L_{\min}\). \(\square\)

The chaining is tight on the \(17\)-cycle: \(u_1=16\), \(u_2=40\), and
\(16^{\delta}/2=40.5\). It is {[}SdW{]} Lemma 6 in the variable \(u\):
their \(x_{i+1}<b^{\delta}x_i^{\delta}\) with
\(b=(1+X_0^{-1})/2^{1/\delta}\) becomes exact here, with no \(X_0\) in
the constant, because \(u=y-1\) is the conjugated variable itself. The
bound on \(o\) is their Lemma 7 read as a lower bound on the least
element rather than as an upper bound on \(\Lambda\); their constant
\(c_m=2^{(m/\delta)(\delta-1)/(\delta^m-1)}\,b^{\delta/(\delta-1)-m/(\delta^m-1)}\)
is, at \(b^{\delta}=1/2\), the \(2^{-(B-m)/((\delta-1)B)}\) of Lemma 3.

\textbf{Lemma 4 (admissible lengths).} If an \(m\)-cycle has
\(x_{\min}\ge X_0\), then \(0<\Lambda<m/(X_0-1)<\log3\), so
\(o=\lceil Kx\rceil\) and \[
\Lambda(K)=\log3\,\bigl(1-\{Kx\}\bigr)<\frac m{X_0-1}.
\] Hence the admissible lengths, those a cycle above the floor can have,
are exactly the \(K\) with \(1-\{Kx\}<m/((X_0-1)\log3)\), and below any
bound they are finite in number.

\emph{Proof.} Lemma 2 gives \(0<\Lambda<m/(x_{\min}-1)\le m/(X_0-1)\),
and \(m/(X_0-1)<\log3\) at every \((m,X_0)\) in Section 5. Then
\(0<o\log3-K\log2<\log3\) puts \(o\) strictly between \(Kx\) and
\(Kx+1\), and \(Kx\) is irrational, so \(o=\lceil Kx\rceil\) and
\(\Lambda=\log3\,(o-Kx)=\log3\,(1-\{Kx\})\). Finiteness below a bound is
immediate. \(\square\)

\emph{How the finite set is listed.} In practice the returns of
\(\{Kx\}\) into the window are walked: from the two least returns
\(q_1,q_2\), found among the convergents and semiconvergents of \(x\),
each further member is the previous one plus \(q_1\), \(q_2\) or
\(q_1+q_2\), which is the classical three-distance structure {[}3D{]};
the least member is Crandall's lower bound on the length. Nothing below
depends on that recursion being the right one. Theorem 7 quantifies over
\emph{every} admissible \(K\) below the ceiling, and the lists used to
verify it were produced twice by unrelated means, by the walk and by an
exhaustive integer sieve over \(K=iQ+j\) at \(2^{256}\) scale, which
agree (Section 8).

{[}SdW{]} bound the same set from both ends before listing it: Lemma 10
and Corollary 11 give the least convergent denominator the floor admits,
Lemma 16 caps the length through the largest partial quotient below that
index, and Section 7 lists what lies between by an approximation
lattice. The enumeration below replaces all three.

\textbf{Lemma 5 (contracting prefixes, and the sieve the floor uses).}
Let \(w\) be the first \(j\) parities of the orbit of \(y\), with \(a\)
of them odd. Then \[
2^j g^j(y)=3^a y-D(w),\qquad D(w)\ge0,
\] with \(D\) a function of the word alone. The parity of \(g^i(y)\)
depends only on \(y\bmod2^{i+1}\), so \(w\) depends only on
\(y\bmod2^j\). Consequently, if \(3^a<2^j\) then \(g^j(y)<y\) for
\textbf{every} \(y\) in that residue class, with no lower threshold on
\(y\), and a descent verification may skip the class entirely.

\emph{Proof.} Induction on \(j\). At \(j=0\), \(D=0\). An even step has
\(2g(y')=y'\), so \(2^{j+1}g^{j+1}(y)=2^jg^j(y)\) and \(D\) is
unchanged; an odd step has \(2g(y')=3y'-1\), so
\(2^{j+1}g^{j+1}(y)=3\bigl(3^ay-D\bigr)-2^j\) and \(D\) becomes
\(3D+2^j\). Both keep \(D\ge0\) and depend on the word alone. Since
\(3^a\) is odd, \(2^ig^i(y)=3^{a_i}y-D_i\) determines \(g^i(y)\bmod2\)
from \(y\bmod2^{i+1}\). Finally \(g^j(y)=(3^ay-D)/2^j\le3^ay/2^j<y\)
when \(3^a<2^j\). \(\square\)

The absence of a threshold is the sign flip again: for \(3n+1\) the word
constant is added, so a contracting prefix drops only the members of the
class above an explicit bound, and the sieve must carry that bound. Here
it is subtracted and the class falls entire. This is what makes the
floor computation of Section 1 exact rather than approximate: at
\(j=24\) the \(286581\) classes that never contract are the only ones
walked, a count that is OEIS A076227(24).

\section{4. Rhin's ceiling}\label{rhins-ceiling}

\textbf{Proposition 6.} Rhin's Proposition {[}R87, p.~160, (7){]} states
that for integers \(u_0,u_1,u_2\) with \(H=\max(|u_1|,|u_2|)\ge2\), \[
|u_0+u_1\log2+u_2\log3|\ \ge\ H^{-13.3},
\] with no further constant. For an \(m\)-cycle of length \(K\ge2\) with
an even step, \((u_0,u_1,u_2)=(0,-K,o)\) and \(H=K\), so \[
2^{L_{\min}(K,m)}\ <\ x_{\min}-1\ <\ \frac m{\Lambda}\ \le\ m\,K^{13.3}.
\] Let \(K_3(m)\) be the least integer such that
\(2^{L_{\min}(K,m)}\ge2m\,K^{13.3}\) for every \(K\ge K_3(m)\). Then
every \(m\)-cycle has \(K<K_3(m)\).

\emph{Proof.} Lemmas 2 and 3 with Rhin's bound. The function
\(K\mapsto L_{\min}(K,m)-13.3\log_2K-\log_2(2m)\) is convex, negative at
\(K=2\), and has a larger root; beyond it the function increases, so
\(K_3(m)\) is well defined, and the doubled constant covers
\(m/\Lambda+1\). \(\square\)

Rhin's bound is the one external input of this note. It is a theorem,
not a hypothesis; it is external in the sense that nothing here
re-proves it and no Lean statement carries it. {[}SdW{]} Lemma 12 states
it as \(\Lambda>e^{-13.3(0.46057+\log K)}\) in their odd count \(K\),
which is (7) at \(H=K+L<1.000001\,\delta K\) with
\(0.46057=\log\delta\); the laboratory's Paper A and its cycle-gap
branch carried that form, weaker by the factor \(\delta^{13.3}=457\)
when \(H\) is the length itself, and it is (7) that is used here. The
general two-logarithm bound of Laurent--Mignotte--Nesterenko {[}LMN{]}
would replace \(13.3\) by an exponent in the thousands at these heights
and is not competitive; Rhin's constant is specific to \(\log2\) and
\(\log3\).

\section{5. The theorem}\label{the-theorem}

\textbf{Theorem 7.} Let \(X_0=2^{51}\). For every \(1\le m\le58\) and
every length \(K<K_3(m)\) with \(\Lambda(K)<m/(X_0-1)\), \[
2^{L_{\min}(K,m)}\ \ge\ \frac m{\Lambda(K)} .
\] Consequently, with Rhin's bound, the \(3n-1\) map has no \(m\)-cycle
with \(1\le m\le58\) whose least element is at least \(2^{51}\), and by
the floor no \(m\)-cycle with \(m\le58\) other than \((5,7,10)\) and the
cycle at \(17\). For \(m\le2\) this is weaker than {[}S07{]}, which
needs no floor; the content of the theorem is the range \(3\le m\le58\).

\emph{Proof.} An \(m\)-cycle above the floor has \(K<K_3(m)\) by
Proposition 6, its length is admissible by Lemma 4, and Lemmas 2 and 3
give \(2^{L_{\min}(K,m)}<x_{\min}-1\le m/\Lambda(K)\), contradicting the
display. The display is a finite exact computation: for each \(m\) the
admissible lengths below \(K_3(m)\) are the finitely many \(K\) of Lemma
4, listed by two unrelated methods that agree, and each is tested at
eighty digits. For \(m\le52\) there is no admissible length below
\(K_3(m)\) at all, so the conclusion there rests on Lemmas 2 and 4 and
Proposition 6 alone. For \(53\le m\le58\) the admissible lengths exist
and fail the display by the margins of Table 1. \(\square\)

\textbf{Table 1.} Rows at the floor \(2^{51}\). ``Margin'' is
\(\max\log_2\bigl(m/(\Lambda(K)\,2^{L_{\min}})\bigr)\) over the
admissible \(K<K_3(m)\), negative when every admissible length is
excluded.

\begingroup\footnotesize

\begin{longtable}[]{@{}
  >{\raggedright\arraybackslash}p{(\linewidth - 10\tabcolsep) * \real{0.1667}}
  >{\raggedright\arraybackslash}p{(\linewidth - 10\tabcolsep) * \real{0.1667}}
  >{\raggedright\arraybackslash}p{(\linewidth - 10\tabcolsep) * \real{0.1667}}
  >{\raggedright\arraybackslash}p{(\linewidth - 10\tabcolsep) * \real{0.1667}}
  >{\raggedright\arraybackslash}p{(\linewidth - 10\tabcolsep) * \real{0.1667}}
  >{\raggedright\arraybackslash}p{(\linewidth - 10\tabcolsep) * \real{0.1667}}@{}}
\toprule\noalign{}
\begin{minipage}[b]{\linewidth}\raggedright
\(m\)
\end{minipage} & \begin{minipage}[b]{\linewidth}\raggedright
\(m/(X_0-1)\)
\end{minipage} & \begin{minipage}[b]{\linewidth}\raggedright
\(K_3(m)\)
\end{minipage} & \begin{minipage}[b]{\linewidth}\raggedright
admissible \(K<K_3\)
\end{minipage} & \begin{minipage}[b]{\linewidth}\raggedright
least admissible
\end{minipage} & \begin{minipage}[b]{\linewidth}\raggedright
margin (bits)
\end{minipage} \\
\midrule\noalign{}
\endhead
\bottomrule\noalign{}
\endfoot
\bottomrule\noalign{}
\endlastfoot
1 & \(4.4\cdot10^{-16}\) & 155 & 0 & --- & --- \\
2 & \(8.9\cdot10^{-16}\) & 495 & 0 & --- & --- \\
5 & \(2.2\cdot10^{-15}\) & 3926 & 0 & --- & --- \\
10 & \(4.4\cdot10^{-15}\) & 57129 & 0 & --- & --- \\
20 & \(8.9\cdot10^{-15}\) & 8393107 & 0 & --- & --- \\
30 & \(1.3\cdot10^{-14}\) & 1094917645 & 0 & --- & --- \\
40 & \(1.8\cdot10^{-14}\) & 134716197551 & 0 & --- & --- \\
50 & \(2.2\cdot10^{-14}\) & 15974496522785 & 0 & --- & --- \\
53 & \(2.4\cdot10^{-14}\) & 66574098182843 & 1 & 64789416887513 &
\(-549.2\) \\
54 & \(2.4\cdot10^{-14}\) & 107084917386112 & 2 & 64789416887513 &
\(-328.1\) \\
55 & \(2.4\cdot10^{-14}\) & 172208666666214 & 4 & 64789416887513 &
\(-188.5\) \\
56 & \(2.5\cdot10^{-14}\) & 276877898691766 & 6 & 64789416887513 &
\(-100.4\) \\
57 & \(2.5\cdot10^{-14}\) & 445072680198484 & 10 & 64789416887513 &
\(-44.8\) \\
58 & \(2.6\cdot10^{-14}\) & 715295665905327 & 16 & 64789416887513 &
\(-9.8\) \\
59 & \(2.6\cdot10^{-14}\) & 1149356600587428 & 27 & 64789416887513 &
\(+12.4\) \\*
\end{longtable}

\endgroup

The closest exclusion, at \(m=58\), has \(9.8\) bits to spare at
\(K=64789416887513\), the least admissible length at this floor; the
same length is the first to survive at \(m=59\). At \(m=59\) the
template leaves two lengths, \[
64789416887513,\quad83130157078217,
\] with \(12.4\) and \(5.0\) bits of room; they are near-convergent
lengths of \(x\) on the expanding side. They are the output of the
template's last step, not its input: what removes each is a higher
floor, at \(2^{51.87}\) and \(2^{55.18}\) respectively (Lemma 2's
\(x_{\min}-1<m/\Lambda\)).

\textbf{Table 2.} What each floor buys, by the same tables.

\begingroup\footnotesize

\begin{longtable}[]{@{}
  >{\raggedright\arraybackslash}p{(\linewidth - 6\tabcolsep) * \real{0.2500}}
  >{\raggedright\arraybackslash}p{(\linewidth - 6\tabcolsep) * \real{0.2500}}
  >{\raggedright\arraybackslash}p{(\linewidth - 6\tabcolsep) * \real{0.2500}}
  >{\raggedright\arraybackslash}p{(\linewidth - 6\tabcolsep) * \real{0.2500}}@{}}
\toprule\noalign{}
\begin{minipage}[b]{\linewidth}\raggedright
floor
\end{minipage} & \begin{minipage}[b]{\linewidth}\raggedright
\(m\) excluded through
\end{minipage} & \begin{minipage}[b]{\linewidth}\raggedright
first open \(m\)
\end{minipage} & \begin{minipage}[b]{\linewidth}\raggedright
lengths left at the first open \(m\)
\end{minipage} \\
\midrule\noalign{}
\endhead
\bottomrule\noalign{}
\endfoot
\bottomrule\noalign{}
\endlastfoot
\(2^{40}\) & 44 & 45 & 72448885240, 82888745831, 93328606422,
103768467013 \\
\(2^{44}\) & 49 & 50 & 539722056247, 757698850864, 975675645481,
1193652440098 \\
\(2^{48}\) & 49 & 50 & 1193652440098 \\
\(2^{49}\) & 54 & 55 & 9767196315401 \\
\(2^{50}\) & 56 & 57 & 28107936506105 \\
\(2^{51}\) & 58 & 59 & 64789416887513, 83130157078217 \\
\(2^{56}\) & 63 & 64 & 766512153894657 \\
\(301\cdot2^{50}\) ({[}SdW{]}'s floor) & 63 & 64 & 766512153894657 \\
\(2^{60}\) & 68 & 69 & 9881527843552324 \\
\(2^{68}\) & 82 & 83 & 6094436882695943503, 6724555128221608268 \\*
\end{longtable}

\endgroup

Two comparisons calibrate the method. First, on the \(3n+1\) side, with
the window on the contracting side of \(x\), the same enumeration and
the same two tests at the floor of {[}SdW{]}, \(301\cdot2^{50}\), return
their Lemma 18: no length for \(64\le m\le68\), and at \(69\le m\le72\)
their five pairs \((K,L)\) at the nine places of their table, with the
floors that remove them to rounding: their first pair
\((5750934602875680,\,3364081086781987)\) falls at \(576.2\), \(584.6\),
\(592.9\), \(601.3\) times \(2^{50}\) for \(m=69,\dots,72\) against
their \(577\), \(585\), \(593\), \(602\), their second at \(623.6\),
\(632.4\) against \(624\), \(633\), and their three at \(m=72\) at
\(308.2\), \(666.8\), \(705.3\) against \(309\), \(667\), \(706\). The
same run lists one point more at \(m=72\), the double of their second
pair, with \(0.9\) bits of room and a removing floor of
\(316.2\cdot2^{50}\); their table does not carry it, and the text does
not say whether non-primitive pairs are discarded. On the \(3n-1\) side
the same floor gives \(m\le63\): the length \(766512153894657\) sits
inside the window at \(m=64\) on the expanding side, where the
contracting side has nothing below the ceiling until \(m=69\). The five
values between \(63\) and \(68\) are the sign of the map, not a missing
stage. Second, Hercher's \(91\) at \(695\cdot2^{60}\) against the plain
template's \(82\) at \(2^{68}\) prices his valley arrangement at about
nine values of \(m\), the one refinement left to transpose. The floor
enters only through the admissibility threshold, in steps: \(2^{52}\) to
\(2^{55}\) buy nothing over \(2^{51}\); \(2^{56}\) buys \(m\le63\),
\(2^{59}\) buys \(64\), \(2^{60}\) buys \(68\).

\section{6. The cycles that exist pass the same
test}\label{the-cycles-that-exist-pass-the-same-test}

The inequalities were checked where they must not exclude. With the
floor set at \(17\) and \(m=2\), the admissible lengths include \(11\)
and it survives the display; with the floor at \(5\) and \(m=1\), length
\(3\) survives; with the floor at \(18\), length \(11\) no longer
appears, since the \(17\)-cycle's least element is below it. The cycle
equation of Lemma 2 is verified in exact rationals on all three cycles,
and Lemma 1's run length and identity on every odd \(y<4000\).

\section{7. What this note does not
do}\label{what-this-note-does-not-do}

It excludes no Juggler cycle. The Juggler's cycle words are, letter for
letter, the words of the negative Collatz cycles ({[}A{]}, Section 5.9),
but a word shape does not transport a realization, so a \(3n-1\)
\(m\)-cycle theorem constrains Juggler cycle words and nothing more. It
does not settle the \(3n-1\) cycle question, \(m\) being bounded, and it
says nothing about divergence on either map. Hercher's refinement is not
transposed; the lattice of admissible periods that Eliahou {[}E93{]}
builds on the positive side is not transposed.

Lemmas 1 and 3 are machine-checked
(\texttt{\nolinkurl{Problems.Collatz.NegativeMCycles}}, twenty
declarations, Lean 4 with Mathlib, no \texttt{\nolinkurl{sorry}} and
nothing off the kernel). The formalization writes the start of an odd
run as \(y=2^a m+1\) with \(m\) odd, which removes every truncated
subtraction and puts the run in closed form,
\(g^k(2^a m+1)=3^k 2^{a-k} m+1\) for \(k\le a\); Lemma 1's four clauses
and Lemma 3's two halves are read off it, and the known cycles are
checked against the same definitions inside the kernel. Formalizing it
sharpened one hypothesis: the closed form needs no parity assumption on
\(m\), since the powers of two alone keep the state odd through the run,
and oddness of \(m\) is used only to know that the run stops at \(a\).
Lemma 2 is not in Lean: its content is the cycle equation and the bound
on \(\Lambda\), which are real-analytic. Its first clause,
\(\Lambda>0\), is the negative-cycle expansion already in Lean on the
conjugate side (\texttt{\nolinkurl{neg_cycle_expanding}}), together with
the negative-side finance and word shape
(\texttt{\nolinkurl{neg_cycle_finance}},
\texttt{\nolinkurl{neg_prefix_noncontracting}}); the conjugation
\(x=-y\) that carries one to the other is stated in this note and is not
itself formalized.

\section{8. Verification}\label{verification}

\texttt{python\ -m\ research.juggler\_sequence.negative\_m\_cycles}
recomputes both tables in a few seconds and writes
\texttt{\nolinkurl{data/research/juggler/negative_m_cycles/summary.json}};
\texttt{\nolinkurl{tests/research/juggler_sequence/test_negative_m_cycles.py}}
checks Lemma 1 on every odd \(y<4000\), the cycle equation on the three
cycles, the tightness of Lemmas 1 and 3 at \(17\), that the known cycles
pass, that \(K_3\) is a ceiling, that the three-gap walk agrees with a
direct scan, that the archived tables reproduce, and that the same
machinery with the window on the contracting side returns {[}SdW{]}
Lemma 18 at their floor. The floor's certificate is with the branch
\texttt{\nolinkurl{negative_floor_3x1}}: the CPU chunk reports below
\(2^{44}\), and above it the GPU sweep's chunk reports, its calibration
on the CPU's range and the three spot-check windows
(\texttt{python\ -m\ research.juggler\_sequence.negative\_floor\_gpu\ calibrate},
\texttt{sweep\ 44\ 51}, \texttt{\nolinkurl{spot}}).
\texttt{lake\ build\ Problems.Collatz.NegativeMCycles} checks Lemmas 1
and 3; the module is in the \texttt{\nolinkurl{Problems}} barrel, so the
default build covers it, and
\texttt{\nolinkurl{tests/research/juggler_sequence/test_negative_m_cycles_lean.py}}
holds it to no \texttt{\nolinkurl{sorry}}, no
\texttt{\nolinkurl{native_decide}} and the three standard axioms.

\begin{verbatim}
Repository:  https://github.com/sneakyweasel/btlab
Commit:      d6e8a4ae20a5c82f46fc5e673a1bf8d1a344e29c
Lean:        leanprover/lean4:v4.33.1
Mathlib:     v4.33.1 (lake-manifest rev 0df444a360eaa60ab8c11dca51a86af692955474)
Build:       lake build Problems.Collatz.NegativeMCycles   (from formal/)
Tables:      python -m research.juggler_sequence.negative_m_cycles
Floor:       python -m research.juggler_sequence.negative_floor_gpu sweep 44 51
Recheck:     python tools/check_3n_minus_1_note_numeric.py
\end{verbatim}

The commit is the repository state that produced the tables, the floor
records and the Lean module; every file those commands read or write is
byte-identical there to the version this paper reports. A later
editorial commit of this text, including the one that adds the build
tooling for this paper, does not change them.

\section{Acknowledgments and use of
AI}\label{acknowledgments-and-use-of-ai}

Large language models assisted the development of this work, including
prose, proposed proof arguments, the Lean formalization, the CUDA
verifier, and the computations. The sources of Simons--de Weger 2005 and
Rhin 1987 were read directly and the constants taken from them, not from
memory; the transposition to the \(3n-1\) side, the enumeration, the
tables and their independent recomputation were produced and checked in
that collaboration. AI assistance is not independent mathematical
validation. The author is responsible for the statements, the proofs,
the code, and the decision to make this version public.

\section{Availability and version}\label{availability-and-version}

This is version 1.0.0 of Paper D, of 21 September 2026, at the floor
\(2^{51}\). It is a preprint and has not been refereed. The manuscript,
the probe that computes the tables, the independent check that
recomputes them, the Lean module, the CUDA verifier and the floor
certificate with its chunk reports are in the repository
\href{https://github.com/sneakyweasel/btlab}{sneakyweasel/btlab} at the
commit named in Section 8; \texttt{python\ tools/build\_paper\_d.py}
rebuilds this document from its canonical Markdown and
\texttt{\nolinkurl{--check}} verifies every generated copy against the
manifest. The verification floor is a computation of this laboratory,
not a result from the literature, and is labelled as one throughout; a
higher floor raises the theorem's \(m\) and would be a new version
rather than a correction of this one.

\section{References}\label{references}

\begingroup\small

\begin{itemize}
\tightlist
\item
  {[}H23{]} C. Hercher, ``There are no Collatz m-cycles with m ≤ 91,''
  \emph{J. Integer Seq.} 26 (2023), Article 23.3.5. arXiv:2201.00406.
\item
  {[}B21{]} D. Barina, ``Convergence verification of the Collatz
  problem,'' \emph{J. Supercomput.} 77 (2021), 2681--2688.
\item
  {[}S07{]} J. L. Simons, ``A simple (inductive) proof for the
  non-existence of 2-cycles of the 3x+1 problem,'' \emph{J. Number
  Theory} 123 (2007), 10--17. Section 6 treats \(3x-1\).
\item
  {[}S08{]} J. L. Simons, ``On the (non-)existence of m-cycles for
  generalized Syracuse sequences,'' \emph{Acta Arith.} 131 (2008),
  217--254.
\item
  {[}SdW{]} J. L. Simons and B. M. M. de Weger, ``Theoretical and
  computational bounds for m-cycles of the 3n+1-problem,'' \emph{Acta
  Arith.} 117 (2005), 51--70.
\item
  {[}Si05{]} J. L. Simons, ``On the nonexistence of 2-cycles for the
  3x+1 problem,'' \emph{Math. Comp.} 74 (2005), 1565--1572.
\item
  {[}R87{]} G. Rhin, ``Approximants de Pad\'{e} et mesures effectives
  d'irrationalit\'{e},'' \emph{S\'{e}minaire de Th\'{e}orie des Nombres,
  Paris 1985--86}, Progress in Mathematics 71, Birkh\"{a}user, 1987,
  155--164.
\item
  {[}3D{]} V. T. S\'{o}s, ``On the distribution mod 1 of the sequence
  \(n\alpha\),'' \emph{Ann. Univ. Sci. Budapest. E\"{o}tv\"{o}s Sect.
  Math.} 1 (1958), 127--134; S. \'{S}wierczkowski, ``On successive
  settings of an arc on the circumference of a circle,'' \emph{Fund.
  Math.} 46 (1959), 187--189; N. B. Slater, ``Gaps and steps for the
  sequence \(n\theta\) mod 1,'' \emph{Proc. Cambridge Philos. Soc.} 63
  (1967), 1115--1123.
\item
  {[}LMN{]} M. Laurent, M. Mignotte and Y. Nesterenko, ``Formes
  lin\'{e}aires en deux logarithmes et d\'{e}terminants
  d'interpolation,'' \emph{J. Number Theory} 55 (1995), 285--321.
\item
  {[}Si03{]} M. K. Sinisalo, ``On the minimal cycle lengths of the
  Collatz sequences,'' preprint, University of Oulu, 2003.
\item
  {[}E93{]} S. Eliahou, ``The 3x+1 problem: new lower bounds on
  nontrivial cycle lengths,'' \emph{Discrete Math.} 118 (1993), 45--56.
\item
  {[}St77{]} R. P. Steiner, ``A theorem on the Syracuse problem,''
  \emph{Proc. 7th Manitoba Conf. Numerical Math.} (1977), 553--559.
\item
  {[}L85{]} J. C. Lagarias, ``The 3x+1 problem and its
  generalizations,'' \emph{Amer. Math. Monthly} 92 (1985), 3--23; and
  the annotated bibliographies, arXiv:math/0309224 and math/0608208.
\item
  {[}A{]} P. Cochin, ``Lower Bounds for Cycle Lengths in the Juggler
  Map'' (Paper A), Zenodo version 1.0.0, 9 September 2026,
  \href{https://doi.org/10.5281/zenodo.22676453}{doi:10.5281/zenodo.22676453};
  concept DOI
  \href{https://doi.org/10.5281/zenodo.22676452}{10.5281/zenodo.22676452}.
\item
  OEIS A037084, comment on the \(3x-1\) verification to \(10^8\).
\end{itemize}

\endgroup
\end{document}
